\documentclass[conference]{IEEEtran}
\usepackage{cite}
\usepackage{amsmath,amssymb,amsfonts}
\usepackage{algorithmic}
\usepackage{algorithm}
\usepackage{graphicx}
\usepackage{textcomp}
\usepackage{xcolor}
\usepackage{booktabs}
\usepackage{multirow}
\usepackage{hyperref}
\usepackage{braket}
\usepackage{float}

\hypersetup{
    colorlinks=true,
    linkcolor=blue,
    citecolor=blue,
    urlcolor=blue
}

\def\BibTeX{{\rm B\kern-.05em{\sc i\kern-.025em b}\kern-.08em
    T\kern-.1667em\lower.7ex\hbox{E}\kern-.125emX}}

\begin{document}

\title{Quantum-Inspired Hybrid Multi-Robot Path Planning\\with QAOA-Based Coordination}

\author{
\IEEEauthorblockN{Author Name}
\IEEEauthorblockA{Department of Computer Science\\
University Name\\
City, Country\\
email@university.edu}
}

\maketitle

\begin{abstract}
This paper presents a hybrid quantum-classical framework for multi-robot path planning in dynamic two-dimensional grid environments. The system addresses two longstanding challenges in autonomous robotics: (1)~single robots becoming trapped in local minima during navigation, and (2)~the exponential growth of computational cost when coordinating multiple robots. Our approach integrates five algorithms in a layered architecture: A* search for global planning, Artificial Potential Fields~(APF) for real-time obstacle avoidance, a novel Quantum-Inspired Q-Learning method for escaping local minima, a Hybrid Planner that unifies these three layers, and the Quantum Approximate Optimization Algorithm~(QAOA) for multi-robot path coordination. The key contributions include a Grover-inspired rotation-based Q-Learning update rule that encodes action values as bounded qubit rotation angles, and a QUBO formulation that enables QAOA to select collision-free path combinations across multiple robots. Experiments conducted on $20 \times 20$ grids with obstacle densities ranging from 10\% to 30\% demonstrate that the quantum-inspired Q-Learning agent successfully escapes U-shaped traps where APF fails, and the QAOA optimizer achieves 98.7\% of brute-force optimal quality while offering polynomial scaling potential for larger robot teams.
\end{abstract}

\begin{IEEEkeywords}
Multi-robot path planning, quantum-inspired computing, Q-learning, QAOA, artificial potential fields, hybrid planning, combinatorial optimization
\end{IEEEkeywords}

% ===================================================================
\section{Introduction}
% ===================================================================

Multi-robot path planning (MRPP) is a fundamental problem in autonomous systems, with direct applications in warehouse logistics, autonomous vehicles, search and rescue operations, and manufacturing automation~\cite{hart1968formal, khatib1986real}. As organizations such as Amazon deploy hundreds of thousands of mobile robots in fulfillment centers, the need for efficient, collision-free path planning algorithms has never been more pressing.

The MRPP problem requires each robot to navigate from its starting position to a designated goal while avoiding both static obstacles (walls, structures) and dynamic obstacles (other robots, moving objects). This seemingly straightforward task gives rise to two fundamental difficulties. First, at the individual level, reactive planning methods such as Artificial Potential Fields are prone to getting trapped in local minima created by concave obstacles. Second, at the coordination level, selecting the optimal combination of paths for $N$ robots with $K$ candidate paths each requires evaluating $\mathcal{O}(K^N)$ combinations---a problem that becomes computationally intractable even for moderate values of $N$.

Classical approaches address these challenges in isolation. A* search~\cite{hart1968formal} guarantees optimal paths for individual robots but cannot adapt to dynamic environments. Potential field methods~\cite{khatib1986real} provide real-time responsiveness but fail in environments with concave obstacles. Reinforcement learning offers adaptability but suffers from slow convergence and unbounded value functions. Brute-force coordination finds globally optimal solutions but scales exponentially.

Recent advances in quantum computing and quantum-inspired algorithms have opened new avenues for addressing these computational bottlenecks. Quantum reinforcement learning~\cite{dong2008quantum} has demonstrated that encoding action values as quantum states can improve exploration and convergence. Meanwhile, the Quantum Approximate Optimization Algorithm (QAOA)~\cite{farhi2014qaoa} has shown promise in solving combinatorial optimization problems---exactly the class of problem that multi-robot coordination represents.

In this paper, we present a hybrid quantum-classical framework that combines the strengths of five distinct algorithms in a unified architecture. Our key contributions are:

\begin{enumerate}
    \item A \textbf{Grover-inspired rotation-based Q-Learning} update rule that maps Q-values to qubit rotation angles $\theta \in [0, \pi]$, ensuring bounded and self-regulating learning dynamics.
    \item An \textbf{extension from single-robot to multi-robot planning} by formulating path coordination as a Quadratic Unconstrained Binary Optimization (QUBO) problem solved via QAOA.
    \item A \textbf{five-layer hybrid architecture} (A* + APF + Quantum Q-Learning + A* replan + QAOA) that combines global optimality, local reactivity, intelligent trap escape, and scalable multi-robot coordination.
    \item A \textbf{complete numpy-based QAOA simulator} that enables statevector simulation of the optimization circuit without requiring quantum hardware.
\end{enumerate}

The remainder of this paper is organized as follows. Section~\ref{sec:related} reviews related work. Section~\ref{sec:problem} formalizes the problem statement. Section~\ref{sec:methodology} presents our proposed algorithms in detail. Section~\ref{sec:architecture} describes the system architecture. Section~\ref{sec:experiments} discusses experimental setup and results. Section~\ref{sec:conclusion} concludes with a summary of findings and directions for future work.

% ===================================================================
\section{Related Work}
\label{sec:related}
% ===================================================================

\subsection{Classical Path Planning}

The A* algorithm, introduced by Hart, Nilsson, and Raphael~\cite{hart1968formal}, remains the gold standard for optimal shortest-path computation on graphs and grids. It uses a heuristic function to guide its search, expanding nodes in order of estimated total cost $f(n) = g(n) + h(n)$. While A* guarantees optimality when the heuristic is admissible, it requires a complete and static map and does not scale well to multi-robot scenarios.

Khatib~\cite{khatib1986real} introduced the Artificial Potential Field (APF) method for real-time obstacle avoidance. In this approach, the goal generates an attractive force and obstacles generate repulsive forces, with the robot moving along the negative gradient of the combined potential. APF provides elegant real-time reactivity but is well known to suffer from local minima, particularly in environments containing U-shaped or concave obstacles.

\subsection{Reinforcement Learning for Path Planning}

Tabular Q-Learning~\cite{watkins1992qlearning} enables robots to learn navigation policies through trial-and-error interaction with the environment. While flexible and capable of handling dynamic environments, standard Q-Learning uses an $\epsilon$-greedy exploration strategy that treats all states uniformly, and its scalar Q-values are unbounded, which can lead to divergence in certain configurations.

\subsection{Quantum-Inspired Reinforcement Learning}

Dong et al.~\cite{dong2008quantum} introduced the concept of quantum reinforcement learning, where action-selection probabilities are derived from quantum state amplitudes. Subsequent work explored quantum-inspired deep reinforcement learning with qubit experience replay for path planning~\cite{qer2020}, and quantum particle swarm optimization for three-dimensional robot path planning~\cite{caaqpso2021}. A hybrid approach combining Fuzzy A*, Quantum Q-Learning, and APF was proposed for single-robot navigation~\cite{hybridpaper2023}, demonstrating that quantum-inspired updates can help robots escape local minima. However, this prior work was limited to single-robot scenarios and did not address multi-robot coordination.

\subsection{Quantum Optimization for Combinatorial Problems}

The Quantum Approximate Optimization Algorithm (QAOA), proposed by Farhi et al.~\cite{farhi2014qaoa}, provides a variational approach to solving combinatorial optimization problems. QAOA encodes the objective function as a cost Hamiltonian and uses alternating layers of cost and mixer unitaries to search the solution space. While QAOA has been applied to problems such as MaxCut and portfolio optimization, its application to multi-robot path coordination via QUBO formulations remains largely unexplored.

% ===================================================================
\section{Problem Statement}
\label{sec:problem}
% ===================================================================

\subsection{Environment Model}

We consider a discrete two-dimensional grid environment $\mathcal{G}$ of size $M \times M$, where each cell $(r, c)$ is either free or occupied by a static obstacle. Additionally, the environment may contain dynamic obstacles that move along predefined trajectories (bouncing off walls or looping around the grid boundaries).

\subsection{Single-Robot Navigation}

Given a robot with start position $s \in \mathcal{G}$ and goal position $g \in \mathcal{G}$, the task is to find a path $\pi = (s = p_0, p_1, \ldots, p_T = g)$ such that every $p_t$ is a free cell and consecutive positions $p_t, p_{t+1}$ are adjacent (including diagonal neighbors). The path cost is defined as:
\begin{equation}
    C(\pi) = \sum_{t=0}^{T-1} d(p_t, p_{t+1})
\end{equation}
where $d(\cdot, \cdot)$ is the Euclidean distance between adjacent cells ($1$ for cardinal moves, $\sqrt{2}$ for diagonal moves).

\subsection{Multi-Robot Coordination}

Given $N$ robots, each with start position $s_i$ and goal position $g_i$, and $K$ candidate paths per robot, the coordination problem is to select one path $\pi_i^{k_i}$ per robot such that:
\begin{equation}
    \min_{\{k_1, \ldots, k_N\}} \sum_{i=1}^{N} C(\pi_i^{k_i}) + \lambda \sum_{i < j} \text{conflicts}(\pi_i^{k_i}, \pi_j^{k_j})
\end{equation}
where $\text{conflicts}(\cdot, \cdot)$ counts the number of timesteps at which two robot paths occupy the same cell. This is a combinatorial optimization problem with $\mathcal{O}(K^N)$ possible selections.

\subsection{Challenges}

Table~\ref{tab:existing_limitations} summarizes the limitations of existing methods that motivate our hybrid approach.

\begin{table}[h]
\centering
\caption{Limitations of Existing Planning Methods}
\label{tab:existing_limitations}
\begin{tabular}{@{}lll@{}}
\toprule
\textbf{Method} & \textbf{Strength} & \textbf{Limitation} \\
\midrule
A* & Optimal path & Static only, single robot \\
APF & Real-time reaction & Local minima (U-traps) \\
Q-Learning & Learns from experience & Unbounded Q-values \\
Greedy Selection & Fast coordination & Ignores inter-robot collisions \\
Brute-Force & Global optimum & $\mathcal{O}(K^N)$ complexity \\
\bottomrule
\end{tabular}
\end{table}

% ===================================================================
\section{Proposed Methodology}
\label{sec:methodology}
% ===================================================================

\subsection{A* Global Planner}

We implement A* search with support for eight-connected movement (four cardinal and four diagonal directions). The heuristic function is the Euclidean distance:
\begin{equation}
    h(n) = \sqrt{(n_r - g_r)^2 + (n_c - g_c)^2}
\end{equation}
which is admissible for eight-connected grids and guarantees optimality. Corner-cutting through diagonal wall gaps is explicitly prevented. The algorithm maintains a priority queue ordered by $f(n) = g(n) + h(n)$, where $g(n)$ is the actual cost from the start to node $n$.

\subsection{Artificial Potential Field Local Planner}

The APF planner computes forces at each step to provide reactive navigation. The attractive potential drawing the robot toward the goal is:
\begin{equation}
    U_{\text{att}}(\mathbf{q}) = \frac{1}{2} k_{\text{att}} \|\mathbf{q} - \mathbf{q}_g\|^2
\end{equation}
where $k_{\text{att}}$ is the attractive gain and $\mathbf{q}_g$ is the goal position. The repulsive potential pushing the robot away from each nearby obstacle $o$ is:
\begin{equation}
    U_{\text{rep}}(\mathbf{q}) = \begin{cases}
    \frac{1}{2} k_{\text{rep}} \left(\frac{1}{d(\mathbf{q}, o)} - \frac{1}{\rho_0}\right)^2 & \text{if } d(\mathbf{q}, o) \leq \rho_0 \\
    0 & \text{otherwise}
    \end{cases}
\end{equation}
where $k_{\text{rep}}$ is the repulsive gain and $\rho_0$ is the influence radius. The robot moves to the neighboring cell with the lowest total potential $U(\mathbf{q}) = U_{\text{att}}(\mathbf{q}) + \sum_o U_{\text{rep}}(\mathbf{q})$.

A stuck detection mechanism monitors the last $W$ positions. If all positions lie within a radius $r_s$ of their centroid, the robot is declared stuck in a local minimum.

\subsection{Classical Q-Learning Baseline}

Standard tabular Q-Learning serves as the reinforcement learning baseline. The state space is the set of grid cells $(r, c)$, and the action space consists of four cardinal directions. The Q-value update follows the Bellman equation:
\begin{equation}
    Q(s, a) \leftarrow Q(s, a) + \alpha \left[r + \gamma \max_{a'} Q(s', a') - Q(s, a)\right]
\end{equation}
where $\alpha = 0.1$ is the learning rate, $\gamma = 0.95$ is the discount factor, and the reward structure is: $+100$ for reaching the goal, $-10$ for hitting a wall, $-3$ for revisiting a cell, and $-1$ for each normal step. An $\epsilon$-greedy exploration schedule starts at $\epsilon = 1.0$ (fully random) and decays to $\epsilon_{\min} = 0.05$ over 500 training episodes.

\subsection{Quantum-Inspired Q-Learning}
\label{sec:quantum_ql}

This is a central contribution of our work. Instead of storing scalar Q-values that can grow unbounded, we encode each state-action pair as a qubit rotation angle $\theta \in [0, \pi]$.

\subsubsection{Qubit State Representation}

Each state-action pair $(s, a)$ is associated with a single qubit prepared in the state:
\begin{equation}
    \ket{\psi} = R_y(\theta)\ket{0} = \cos\frac{\theta}{2}\ket{0} + \sin\frac{\theta}{2}\ket{1}
\end{equation}
The action-selection probability is the measurement probability of $\ket{1}$:
\begin{equation}
    P(a|s) = \sin^2\!\left(\frac{\theta_{s,a}}{2}\right)
\end{equation}
This mapping is physically meaningful---it corresponds to the probability of measuring $\ket{1}$ when a qubit is rotated from $\ket{0}$ by angle $\theta$ using an $R_y$ gate. The probability is inherently bounded in $[0, 1]$, unlike scalar Q-values which can diverge to $\pm\infty$.

\subsubsection{Grover-Inspired Update Rule}

Our update rule draws inspiration from Grover's search algorithm, where the amplitude of the target state is amplified through iterative rotations. Given the temporal-difference (TD) error, we compute:
\begin{equation}
    \delta_{\text{TD}} = \underbrace{\text{clip}\!\left(\frac{r}{100} + \gamma \max_{a'} P(a'|s'), \; 0, \; 1\right)}_{\text{TD target (clipped probability)}} - P(a|s)
\end{equation}

The rotation update incorporates self-regulating damping that prevents saturation at the boundaries:
\begin{equation}
    \Delta\theta = \alpha \cdot \delta_{\text{TD}} \cdot D(\theta)
\end{equation}
where the damping factor is:
\begin{equation}
    D(\theta) = \begin{cases}
    \frac{\pi - \theta}{\pi} & \text{if } \delta_{\text{TD}} > 0 \quad \text{(rotate toward } \pi\text{)} \\[4pt]
    \frac{\theta}{\pi} & \text{if } \delta_{\text{TD}} < 0 \quad \text{(rotate toward } 0\text{)}
    \end{cases}
\end{equation}

The angle is then updated and clamped:
\begin{equation}
    \theta_{s,a} \leftarrow \text{clip}\!\left(\theta_{s,a} + \Delta\theta, \; 0.01, \; \pi - 0.01\right)
\end{equation}

This self-damping mechanism ensures that as $\theta$ approaches $\pi$ (high probability), positive updates become progressively smaller, and as $\theta$ approaches $0$ (low probability), negative updates diminish. This mimics the natural behavior of Grover rotations near the target amplitude and prevents the oscillation and divergence problems common in standard Q-Learning.

\subsubsection{Advantages Over Standard Q-Learning}

Table~\ref{tab:ql_comparison} summarizes the key differences between the standard and quantum-inspired approaches.

\begin{table}[h]
\centering
\caption{Standard vs.\ Quantum-Inspired Q-Learning}
\label{tab:ql_comparison}
\begin{tabular}{@{}lll@{}}
\toprule
\textbf{Property} & \textbf{Standard QL} & \textbf{Quantum QL} \\
\midrule
Value range & $(-\infty, +\infty)$ & $[0, 1]$ \\
Update type & Additive scalar & Rotational (Bloch sphere) \\
Exploration & $\epsilon$-greedy & Natural via $\sin^2$ \\
Convergence & Can oscillate/diverge & Self-regulating \\
Physical meaning & None & Qubit measurement prob. \\
\bottomrule
\end{tabular}
\end{table}

\subsection{Hybrid Three-Layer Planner}

The Hybrid Planner integrates A*, APF, and Quantum Q-Learning into a unified decision loop, as described in Algorithm~\ref{alg:hybrid}.

\begin{algorithm}[h]
\caption{Hybrid Three-Layer Path Planner}
\label{alg:hybrid}
\begin{algorithmic}[1]
\REQUIRE Grid $\mathcal{G}$, start $s$, goal $g$, trained angles $\theta$
\STATE $\text{path} \leftarrow$ A*($\mathcal{G}$, $s$, $g$) \COMMENT{Layer 1: Global plan}
\STATE $\text{mode} \leftarrow$ APF
\STATE $\text{current} \leftarrow s$
\WHILE{current $\neq g$ \AND steps $< $ max\_steps}
    \IF{mode = APF}
        \STATE Move via APF toward next waypoint on A* path
        \IF{stuck detected (last $W$ positions within radius $r_s$)}
            \STATE mode $\leftarrow$ QUANTUM
        \ENDIF
    \ELSIF{mode = QUANTUM}
        \STATE Follow quantum policy: $a^* = \arg\max_a \sin^2(\theta_{s,a}/2)$
        \STATE Move to $\text{current} + a^*$
        \IF{escape steps exhausted}
            \STATE path $\leftarrow$ A*($\mathcal{G}$, current, $g$) \COMMENT{Replan}
            \STATE mode $\leftarrow$ APF
        \ENDIF
    \ENDIF
\ENDWHILE
\RETURN path from $s$ to $g$
\end{algorithmic}
\end{algorithm}

The stuck detection mechanism examines a sliding window of the last $W = 20$ positions. If all positions fall within a radius of $r_s = 3.0$ cells from their centroid, the robot is declared stuck and control transfers to the quantum escape layer. After a fixed number of escape steps (default 20), A* replans from the new position and the system returns to APF-guided navigation.

\subsection{QAOA Multi-Robot Path Optimizer}

For multi-robot coordination, we formulate the path selection problem as a QUBO and solve it using QAOA.

\subsubsection{QUBO Formulation}

For $N$ robots with $K$ candidate paths each, we define binary decision variables $x_{i,k} \in \{0, 1\}$, where $x_{i,k} = 1$ indicates that robot $i$ selects candidate path $k$. The total number of binary variables is $n = N \times K$. The objective function to minimize is:
\begin{equation}
\begin{split}
    H_{\text{QUBO}} = & \sum_{i,k} c_{i,k} \cdot x_{i,k} \\
    & + \lambda_c \sum_{\substack{(i,k),(j,l) \\ i \neq j}} \text{conflicts}(i,k,j,l) \cdot x_{i,k} \cdot x_{j,l} \\
    & + \lambda_h \sum_{i} \left(\sum_{k} x_{i,k} - 1\right)^2
\end{split}
\end{equation}

The three terms encode: (1)~individual path costs as linear terms on the diagonal, (2)~inter-robot collision penalties as quadratic terms with weight $\lambda_c = 10.0$, and (3)~one-hot constraints ensuring each robot selects exactly one path, with penalty weight $\lambda_h = 50.0$.

The one-hot constraint penalty expands as:
\begin{equation}
    \left(\sum_k x_{i,k} - 1\right)^2 = -\sum_k x_{i,k} + 2\sum_{k < l} x_{i,k} \cdot x_{i,l} + 1
\end{equation}
using the identity $x^2 = x$ for binary variables.

\subsubsection{Ising Conversion}

The QUBO is converted to an Ising Hamiltonian via the substitution $x = (1 - z)/2$, mapping $x \in \{0, 1\}$ to $z \in \{-1, +1\}$:
\begin{equation}
    H_{\text{Ising}} = \sum_{i<j} J_{ij} z_i z_j + \sum_i h_i z_i + c
\end{equation}
where the coupling matrix $J$, local fields $h$, and constant offset $c$ are derived from the symmetric QUBO matrix.

\subsubsection{QAOA Circuit}

The QAOA circuit consists of $p$ alternating layers of cost and mixer unitaries, starting from the uniform superposition $\ket{+}^{\otimes n}$:
\begin{equation}
    \ket{\psi(\boldsymbol{\gamma}, \boldsymbol{\beta})} = \prod_{\ell=1}^{p} U_M(\beta_\ell) \, U_C(\gamma_\ell) \ket{+}^{\otimes n}
\end{equation}
where:
\begin{itemize}
    \item \textbf{Cost unitary:} $U_C(\gamma) = e^{-i\gamma H_C}$ applies a diagonal phase rotation based on the cost Hamiltonian.
    \item \textbf{Mixer unitary:} $U_M(\beta) = \prod_j e^{-i\beta X_j}$ applies single-qubit $X$-rotations to explore the solution space.
\end{itemize}

The variational parameters $\boldsymbol{\gamma} = (\gamma_1, \ldots, \gamma_p)$ and $\boldsymbol{\beta} = (\beta_1, \ldots, \beta_p)$ are optimized classically using the COBYLA optimizer to minimize the expected cost $\langle \psi | H_C | \psi \rangle$. Multiple random restarts (default 5) help avoid poor local optima in the parameter landscape.

For $N = 3$ robots with $K = 4$ candidate paths, QAOA uses $n = 12$ qubits (after deduplication, 11 effective qubits) to simultaneously explore $2^{11} = 2048$ candidate combinations through quantum interference, compared to $K^N = 64$ combinations for brute-force.

% ===================================================================
\section{System Architecture}
\label{sec:architecture}
% ===================================================================

The system is implemented as a modular Python framework consisting of 13 modules and approximately 7,000 lines of code. Fig.~\ref{fig:architecture} illustrates the layered architecture and data flow.

\begin{figure}[h]
\centering
\fbox{\parbox{0.44\textwidth}{\centering
\small
\textbf{Layer 1: A* (Global Planner)}\\
$\downarrow$ shortest path\\[3pt]
\textbf{Layer 2: APF (Local Reactive Planner)}\\
follows A* route, dodges dynamic obstacles\\[3pt]
$\downarrow$ stuck detected?\\[3pt]
\textbf{Layer 3: Quantum QL (Escape Planner)}\\
trained policy breaks free from local minimum\\[3pt]
$\downarrow$ escaped\\[3pt]
\textbf{Layer 1: A* (Replan from new position)}\\[3pt]
$\downarrow$ candidate paths generated per robot\\[3pt]
\textbf{Layer 4: QAOA (Multi-Robot Coordinator)}\\
selects collision-free path combination
}}
\caption{Five-layer hybrid planning architecture.}
\label{fig:architecture}
\end{figure}

The key modules and their responsibilities are:

\begin{itemize}
    \item \texttt{grid.py}: Grid environment with static and dynamic obstacle support, multi-robot configuration, and varied obstacle shape generation.
    \item \texttt{a\_star.py}: A* search with eight-connected movement, Euclidean heuristic, and expansion order tracking for visualization.
    \item \texttt{apf.py}: Artificial Potential Field planner with parabolic attractive potential, inverse-distance repulsive potential, and stuck detection.
    \item \texttt{q\_learning.py}: Tabular Q-Learning baseline with $\epsilon$-greedy exploration and Q-table visualization support.
    \item \texttt{quantum\_q\_learning.py}: Quantum-inspired Q-Learning with rotation-angle encoding, Grover-inspired updates, and optional Qiskit circuit verification.
    \item \texttt{hybrid\_planner.py}: Three-layer hybrid combining A*, APF, and Quantum QL with automatic mode switching.
    \item \texttt{qaoa\_optimizer.py}: QUBO formulation, Ising conversion, numpy-based statevector QAOA simulator, COBYLA optimization with multi-restart, and optional Qiskit backend.
    \item \texttt{multi\_robot.py}: Candidate path generation using A* with penalty zones for diversity, conflict matrix computation, and classical selection baselines.
    \item \texttt{dynamic\_env.py}: Time-stepped simulation engine with collision detection for robot-robot and robot-obstacle interactions.
\end{itemize}

The technology stack includes Python 3.10+, NumPy for numerical computation, SciPy (COBYLA) for QAOA parameter optimization, Matplotlib for visualization, and PyGame for the interactive real-time simulator. Qiskit and Qiskit Aer serve as optional backends for quantum circuit verification, with the system falling back to pure NumPy simulation when these packages are unavailable.

% ===================================================================
\section{Experiments and Results}
\label{sec:experiments}
% ===================================================================

\subsection{Experimental Setup}

All experiments were conducted on $20 \times 20$ grid environments using Python 3.10 with NumPy as the computational backend. Obstacle environments were generated with varied realistic shapes including rectangular blocks, linear walls, and L-shaped structures at five density levels: 10\%, 15\%, 20\%, 25\%, and 30\%. Reinforcement learning agents (both classical and quantum Q-Learning) were trained for 500 episodes with 300 maximum steps per episode.

\subsection{Single-Robot Trap Escape}

The U-shaped trap experiment evaluates each planner's ability to navigate around a concave obstacle. A U-shaped wall is placed between the robot's start and goal positions, creating a configuration where the attractive force from the goal pulls the robot directly into the trap opening, where repulsive forces from the walls create a local minimum.

Table~\ref{tab:trap_results} presents the results.

\begin{table}[h]
\centering
\caption{U-Shaped Trap Escape Results ($20 \times 20$ Grid)}
\label{tab:trap_results}
\begin{tabular}{@{}lccc@{}}
\toprule
\textbf{Planner} & \textbf{Success} & \textbf{Path Cost} & \textbf{Notes} \\
\midrule
A* & Yes & 15.66 & Optimal (full map) \\
APF & \textbf{No (Stuck)} & --- & Oscillates at opening \\
Classical QL & Yes & 18.00 & 500 episodes training \\
Quantum QL & Yes & 24.00 & Rotation-based policy \\
Hybrid & Yes & 25.07 & A* + APF + Quantum \\
\bottomrule
\end{tabular}
\end{table}

A* finds the optimal path since it has complete map knowledge and can plan around the trap from the outset. APF fails entirely, oscillating at the trap opening where attractive and repulsive forces cancel. Both Q-Learning variants learn to escape by discovering the route around the U-shape during training. The Hybrid planner initially follows APF into the trap, detects the stuck condition, switches to the quantum escape policy, breaks free, and replans with A* from the new position. While its path cost is higher due to the initial trap entry and escape overhead, it reliably reaches the goal.

\subsection{Performance Across Obstacle Densities}

Table~\ref{tab:density_results} presents benchmark results across five obstacle density levels with varied obstacle shapes.

\begin{table}[h]
\centering
\caption{Planner Performance Across Obstacle Densities}
\label{tab:density_results}
\resizebox{\columnwidth}{!}{%
\begin{tabular}{@{}clcccc@{}}
\toprule
\textbf{Density} & \textbf{Planner} & \textbf{Success} & \textbf{Cost} & \textbf{Steps} & \textbf{Time (ms)} \\
\midrule
\multirow{5}{*}{10\%}
 & A* & Yes & 28.0 & 22 & 1.1 \\
 & APF & Yes & 28.9 & 22 & 2.1 \\
 & Classical QL & Yes & 40.0 & 41 & 677 \\
 & Quantum QL & Yes & 42.0 & 43 & 2415 \\
 & Hybrid & Yes & 30.6 & 25 & 2476 \\
\midrule
\multirow{5}{*}{20\%}
 & A* & Yes & 31.6 & 28 & 2.8 \\
 & APF & No & --- & 39 & 8.2 \\
 & Classical QL & Yes & 40.0 & 41 & 604 \\
 & Quantum QL & Yes & 50.0 & 51 & 2846 \\
 & Hybrid & Yes & 62.7 & 60 & 2770 \\
\midrule
\multirow{5}{*}{30\%}
 & A* & Yes & 32.7 & 30 & 1.6 \\
 & APF & No & --- & 35 & 5.3 \\
 & Classical QL & Yes & 38.0 & 39 & 543 \\
 & Quantum QL & No & --- & 501 & 3278 \\
 & Hybrid & No & --- & 479 & 3343 \\
\bottomrule
\end{tabular}%
}
\end{table}

Several trends emerge from these results:

\begin{itemize}
    \item \textbf{A*} consistently finds optimal paths across all densities with the lowest computation time, but it operates on a static snapshot and cannot react to dynamic obstacles.
    \item \textbf{APF} performs well at low densities but fails at 15\% and above due to increased frequency of concave trap configurations.
    \item \textbf{Classical Q-Learning} shows reliable success across densities due to its ability to learn from exploration, though its paths are 30--50\% longer than A*.
    \item \textbf{Quantum Q-Learning} succeeds at low-to-moderate densities but struggles at 25\%+ where the bounded rotation angles may need more training episodes to find viable paths in cluttered environments.
    \item \textbf{The Hybrid planner} remains reliable up to 25\% density by leveraging A* for global planning, APF for local navigation, and the quantum escape layer for trap recovery.
\end{itemize}

\subsection{Q-Value Stability Analysis}

A critical advantage of Quantum Q-Learning over its classical counterpart is the bounded nature of its value representation. During training, classical Q-values were observed to range from approximately $-28$ to $+100$, with the potential for further growth in longer training runs. In contrast, quantum rotation angles remained strictly within $[0.01, 3.13]$ (approximately $[0, \pi]$), demonstrating the self-regulating property of the Grover-inspired update rule.

The damping factor $D(\theta)$ provides a principled mechanism for this stability. As a well-learned action approaches $\theta \to \pi$ (high probability), the damping factor $(\pi - \theta)/\pi \to 0$ prevents further positive rotation. Conversely, as a poor action approaches $\theta \to 0$ (low probability), the damping factor $\theta/\pi \to 0$ prevents negative rotation below zero. This behavior mirrors the amplitude amplification in Grover's algorithm, where the rotation magnitude naturally decreases as the target state's amplitude approaches 1.

\subsection{Multi-Robot QAOA Coordination}

The QAOA optimizer was evaluated on a three-robot environment with a $20 \times 20$ grid, four candidate paths per robot, and four dynamic obstacles. Table~\ref{tab:qaoa_results} compares QAOA with classical selection methods.

\begin{table}[h]
\centering
\caption{Multi-Robot Path Selection (3 Robots, 4 Paths Each)}
\label{tab:qaoa_results}
\begin{tabular}{@{}lcccc@{}}
\toprule
\textbf{Selector} & \textbf{Cost} & \textbf{Conflicts} & \textbf{Score} & \textbf{Scaling} \\
\midrule
Greedy & 90.6 & 0 & 90.6 & $\mathcal{O}(NK)$ \\
Random & 91.7 & 0 & 91.7 & $\mathcal{O}(NK)$ \\
Brute-Force & 90.6 & 0 & 90.6 & $\mathcal{O}(K^N)$ \\
\textbf{QAOA} & \textbf{91.7} & \textbf{0} & \textbf{91.7} & \textbf{Poly.\ (quantum)} \\
\bottomrule
\end{tabular}
\end{table}

QAOA achieves a quality ratio of $90.6/91.7 = 98.7\%$ of brute-force optimal, with zero collisions. On this small instance ($K^N = 64$ combinations), brute-force is computationally faster since the QAOA simulator carries overhead from statevector manipulation and classical parameter optimization.

\subsection{Scalability Analysis}

The true advantage of QAOA emerges at scale. Table~\ref{tab:scalability} illustrates how the number of candidate combinations grows with robot count.

\begin{table}[h]
\centering
\caption{Scaling Comparison: Brute-Force vs.\ QAOA}
\label{tab:scalability}
\begin{tabular}{@{}ccll@{}}
\toprule
\textbf{Robots} & \textbf{Combinations} & \textbf{Brute-Force} & \textbf{QAOA} \\
\midrule
3 & 64 & Feasible & Feasible \\
5 & 1,024 & Feasible & Feasible \\
10 & $1.05 \times 10^6$ & Slow & Feasible \\
20 & $1.10 \times 10^{12}$ & Infeasible & Feasible \\
\bottomrule
\end{tabular}
\end{table}

While brute-force exhaustive search grows exponentially as $\mathcal{O}(K^N)$, QAOA's circuit depth and parameter count grow polynomially in $N$. Our numpy-based statevector simulator is practical for up to 16 qubits ($2^{16} = 65{,}536$ amplitudes), corresponding to 4 robots with 4 paths each. Scaling beyond this on classical hardware would require tensor-network simulation methods, while execution on actual quantum hardware would directly leverage the polynomial scaling advantage.

% ===================================================================
\section{Discussion}
\label{sec:discussion}
% ===================================================================

\subsection{When Each Planner Excels}

Our experiments reveal that no single algorithm dominates across all scenarios. Table~\ref{tab:best_planner} summarizes when each approach is the best choice.

\begin{table}[h]
\centering
\caption{Recommended Planner by Scenario}
\label{tab:best_planner}
\begin{tabular}{@{}ll@{}}
\toprule
\textbf{Scenario} & \textbf{Best Planner} \\
\midrule
Static grid, single robot & A* \\
U-shaped trap / concave obstacles & Quantum QL or Hybrid \\
Dynamic obstacles & Hybrid \\
Multi-robot coordination & QAOA \\
High density ($\geq$30\%) & A* or Classical QL \\
\bottomrule
\end{tabular}
\end{table}

This result underscores the value of the hybrid architecture: by combining complementary algorithms, the system adapts to diverse environmental conditions that would cause any single method to fail.

\subsection{Novelty and Contributions}

Our work extends the state of the art in several ways relative to existing literature:

\begin{enumerate}
    \item \textbf{Rotation-based Q-Learning for grid navigation}: While prior work~\cite{dong2008quantum} explored quantum-inspired RL, our specific formulation combining rotation-angle Q-values with a self-damping update rule for grid-based path planning is novel. The damping factor directly addresses the Q-value divergence problem in a principled manner.

    \item \textbf{Single-to-multi-robot extension}: The closest prior work~\cite{hybridpaper2023} applied a similar hybrid approach (A* + APF + Quantum QL) but only for single robots. We extend this to multiple robots by adding candidate path generation with diversity penalties, pairwise conflict matrix computation, and QAOA-based coordination.

    \item \textbf{QUBO formulation for path coordination}: Formulating multi-robot path selection as a QUBO problem with path costs, collision penalties, and one-hot constraints is a novel application of quantum optimization to grid-based multi-robot coordination.
\end{enumerate}

\subsection{Limitations}

Several limitations should be acknowledged:

\begin{enumerate}
    \item \textbf{Simulation overhead}: The numpy statevector simulator has $\mathcal{O}(2^n)$ memory complexity. For instances exceeding 16 qubits, real quantum hardware or advanced classical simulation techniques (such as tensor networks) would be necessary.

    \item \textbf{Small-scale advantage gap}: At three robots, brute-force is faster than QAOA simulation. The quantum computational advantage materializes at 10+ robots where brute-force becomes infeasible.

    \item \textbf{High-density performance}: The quantum Q-Learning agent struggles at 30\% obstacle density. Longer training (beyond 500 episodes) and adaptive learning-rate schedules may improve performance in highly cluttered environments.

    \item \textbf{Approximation ratio}: The current QAOA implementation achieves approximately 98.7\% of optimal quality. Higher circuit depth ($p > 2$) and more optimization restarts could improve this ratio at the cost of increased computation time.
\end{enumerate}

% ===================================================================
\section{Conclusion and Future Work}
\label{sec:conclusion}
% ===================================================================

We have presented a hybrid quantum-classical framework for multi-robot path planning that addresses the limitations of individual classical methods through a layered architecture. The framework integrates A* for global optimality, APF for local reactivity, quantum-inspired Q-Learning for trap escape, and QAOA for multi-robot coordination. Our experiments demonstrate that the quantum Q-Learning agent with its Grover-inspired rotation-based update rule provides bounded, stable learning dynamics and successfully escapes local minima that trap conventional APF planners. The QAOA optimizer achieves near-optimal (98.7\%) collision-free path selection while offering polynomial scaling potential for larger robot teams.

The complete system, comprising approximately 7,000 lines of Python across 13 modules, runs entirely on classical hardware using NumPy-based quantum simulation, making it accessible without quantum computing infrastructure.

Future work includes:
\begin{itemize}
    \item Deploying the QAOA circuit on real quantum hardware (e.g., IBM Quantum) via Qiskit backends to evaluate performance with actual quantum effects.
    \item Scaling to 10+ robots using tensor-network QAOA simulation on classical hardware.
    \item Developing time-dependent QAOA formulations that account for dynamic obstacles during multi-robot coordination.
    \item Integrating deep quantum reinforcement learning for environments with continuous state spaces.
    \item Exploring warm-start strategies for QAOA parameter initialization to reduce optimization overhead.
\end{itemize}

% ===================================================================
% REFERENCES
% ===================================================================

\begin{thebibliography}{00}

\bibitem{hart1968formal}
P.~E. Hart, N.~J. Nilsson, and B.~Raphael, ``A formal basis for the heuristic determination of minimum cost paths,'' \textit{IEEE Transactions on Systems Science and Cybernetics}, vol.~4, no.~2, pp.~100--107, 1968.

\bibitem{khatib1986real}
O.~Khatib, ``Real-time obstacle avoidance for manipulators and mobile robots,'' \textit{The International Journal of Robotics Research}, vol.~5, no.~1, pp.~90--98, 1986.

\bibitem{watkins1992qlearning}
C.~J.~C.~H. Watkins and P.~Dayan, ``Q-learning,'' \textit{Machine Learning}, vol.~8, no.~3--4, pp.~279--292, 1992.

\bibitem{dong2008quantum}
D.~Dong, C.~Chen, H.~Li, and T.-J.~Tarn, ``Quantum reinforcement learning,'' \textit{IEEE Transactions on Systems, Man, and Cybernetics, Part B (Cybernetics)}, vol.~38, no.~5, pp.~1207--1220, 2008.

\bibitem{farhi2014qaoa}
E.~Farhi, J.~Goldstone, and S.~Gutmann, ``A quantum approximate optimization algorithm,'' \textit{arXiv preprint arXiv:1411.4028}, 2014.

\bibitem{qer2020}
Quantum-inspired deep reinforcement learning with qubit experience replay for robot path planning, \textit{Reference Paper 1 (QER-LPD3QN)}, 2020.

\bibitem{caaqpso2021}
Quantum particle swarm optimization for 3D robot path planning (CAA*QPSO), \textit{Reference Paper 2}, 2021.

\bibitem{hybridpaper2023}
Fuzzy A* with quantum Q-learning and artificial potential fields for single-robot hybrid path planning, \textit{Reference Paper 3}, 2023.

\bibitem{grover1996fast}
L.~K. Grover, ``A fast quantum mechanical algorithm for database search,'' in \textit{Proc.\ 28th Annual ACM Symposium on Theory of Computing (STOC)}, pp.~212--219, 1996.

\bibitem{lucas2014ising}
A.~Lucas, ``Ising formulations of many NP problems,'' \textit{Frontiers in Physics}, vol.~2, p.~5, 2014.

\end{thebibliography}

\end{document}
